%========================================================
% CHAPTER 2 — Fiber-Channel Modeling
%========================================================

The mathematical framework introduced in
Chapter~\ref{ch:01_mathematical_and_physical_framework}
provides the tools required to describe the evolution and characterization of
bipartite polarization states. The purpose of the present chapter is to apply
these tools to the progressive construction of models for the propagation of
polarization-entangled photon pairs through optical fibers.

The modeling strategy proceeds from analytically controlled effective
descriptions toward a more physically accurate frequency-dependent representation of fiber
propagation. The first models retain only the polarization degree of freedom and
separately describe coherent phase evolution, statistical phase fluctuations,
and depolarization. These mechanisms are then combined into an effective
polarization channel for the two-arm fiber configuration.

The monochromatic approximation is subsequently relaxed by introducing the
spectral structure of the photons and the frequency dependence of birefringent
fiber propagation. This leads to a polarization-mode-dispersion (PMD) model in
which the complete polarization--frequency evolution remains unitary for an ideal
lossless fiber, while an effectively mixed polarization state can emerge after
the unresolved spectral degree of freedom is traced out.

Figure~\ref{fig:fiber_modeling_hierarchy} summarizes the modeling hierarchy
adopted throughout the chapter. The different levels represent increasing
modeling scope rather than a strict physical derivation from one model to the
next.

\begin{figure}[htbp]
    \centering

    \begin{tikzpicture}[
        >=Latex,
        every node/.style={font=\small},
        model/.style={
            draw,
            rounded corners,
            align=center,
            text width=0.63\textwidth,
            minimum height=0.72cm,
            inner sep=4pt
        }
    ]

        \node[model] (phase) at (0,0)
        {Fixed coherent relative-phase evolution};

        \node[model] (random) at (0,-1.3)
        {Statistical phase fluctuations and effective dephasing};

        \node[model] (global) at (0,-2.6)
        {Global depolarizing reference model};

        \node[model] (local) at (0,-3.9)
        {Independent local depolarization in arms \(A\) and \(B\)};

        \node[model] (composite) at (0,-5.2)
        {Composite effective polarization channel};

        \node[model] (pmd) at (0,-6.5)
        {Frequency-dependent birefringent propagation and PMD};

        \draw[->] (phase) -- (random);
        \draw[->] (random) -- (global);
        \draw[->] (global) -- (local);
        \draw[->] (local) -- (composite);
        \draw[->] (composite) -- (pmd);

    \end{tikzpicture}

    \caption{
        Modeling hierarchy adopted in this work. The first levels provide
        effective polarization-only descriptions of the relevant propagation
        mechanisms, while the final level explicitly includes the spectral
        dependence of birefringent fiber propagation.
    }
    \label{fig:fiber_modeling_hierarchy}
\end{figure}

%========================================================
% SECTION 2.1 — Residual Phase Model
%========================================================

\section{Residual Phase Model}
\label{sec:reduced_phase_model}

A general lossless optical fiber may produce an arbitrary unitary transformation
of the polarization state. The first propagation model considered in this work is
a reduced description of this general evolution and is motivated by the
polarization-compensation procedure employed in the experimental setup.

Polarization compensation is used to realign the dominant output polarization
basis with the reference basis of the measurement system. Such compensation does
not, however, imply a complete reconstruction or inversion of the full fiber Jones
operator\footnote{The Jones matrix is the \(2\times2\) complex matrix describing
the transformation of the polarization amplitudes through an optical element or
fiber.}. In particular, a residual rotation around the stabilized polarization
axis may remain.

A limitation of this compensation procedure arises when the polarization state is
monitored primarily through measured populations. Since these measurements provide
information on the polarization populations but not directly on their relative
phase, a residual phase \(e^{i\theta}\) may remain unresolved.

A more complete compensation may be achieved by extending the monitoring system
to additional polarization bases, since phase-sensitive measurements can provide
information on residual transformations that cannot be inferred from
computational-basis populations alone.

Within the computational basis adopted in this thesis, the stabilized polarization
axis is identified with the \(z\)-axis of the Bloch sphere. The remaining coherent
transformation can therefore be represented by a relative phase between the
horizontal and vertical polarization components, reducing the general polarization
transformation to a single phase parameter.

For one polarization qubit, the corresponding unitary operator is

\begin{equation}
U(\theta)
=
\begin{pmatrix}
1 & 0\\
0 & e^{i\theta}
\end{pmatrix},
\label{eq:single_qubit_residual_phase_unitary}
\end{equation}

defined up to an irrelevant global phase. Equivalently, this transformation
corresponds, up to a global phase factor, to a rotation around the \(z\)-axis of
the Bloch sphere.

At this modeling level, the phase is assumed to be fixed for each propagation
realization. Statistical fluctuations of the same residual phase parameter are
introduced separately in the following model.

If the two fiber arms introduce residual phases \( \theta_A \) and \( \theta_B \),
respectively, their joint action is

\begin{equation}
U_A(\theta_A)\otimes U_B(\theta_B).
\label{eq:two_arm_residual_phase_unitary}
\end{equation}

Applying this transformation to the reference Bell state gives

\begin{equation}
|\psi(\theta_A,\theta_B)\rangle
=
\frac{1}{\sqrt{2}}
\left(
e^{i\theta_B}|01\rangle
+
e^{i\theta_A}|10\rangle
\right).
\label{eq:two_arm_residual_phase_bell_state}
\end{equation}

Factoring out the global phase \(e^{i\theta_B}\) yields

\begin{equation}
|\psi(\theta)\rangle
=
\frac{1}{\sqrt{2}}
\left(
|01\rangle
+
e^{i\theta}|10\rangle
\right),
\qquad
\theta=\theta_A-\theta_B.
\label{eq:residual_phase_bell_state}
\end{equation}

The two independent local phase parameters can therefore be reduced to the single
effective parameter \(\theta\). This relative phase does not modify the
computational-basis populations, but it affects phase-sensitive polarization
correlations.

With the convention adopted throughout this work, the same state may be generated
by assigning the complete relative transformation to arm \(A\),

\begin{equation}
U_\theta
=
U_A(\theta)\otimes I_B.
\label{eq:residual_phase_single_arm_convention}
\end{equation}

This convention simplifies the analytical and numerical representation without
changing any observable prediction. The parameter \(\theta\) represents the
effective differential residual polarization phase between the two arms.

The term \emph{residual phase model} will be used throughout this work for this
one-parameter effective description. Its relation to the more general
fixed-frequency birefringent transformation is derived explicitly in
Section~\ref{subsec:fixed_frequency_birefringence_residual_phase}.

%--------------------------------------------------------
% Subsection 2.1.1 — Correlation Structure
%--------------------------------------------------------

\subsection{Correlation Structure}
\label{subsec:residual_phase_correlation_structure}

For the state in Eq.~\eqref{eq:residual_phase_bell_state}, the correlation tensor
defined in Chapter~\ref{ch:01_mathematical_and_physical_framework} is

\begin{equation}
T(\theta)
=
\begin{pmatrix}
\cos\theta & -\sin\theta & 0\\
\sin\theta & \cos\theta & 0\\
0 & 0 & -1
\end{pmatrix}.
\label{eq:residual_phase_correlation_tensor}
\end{equation}

The phase therefore rotates the transverse correlation structure in the
\(xy\)-plane while leaving the longitudinal anticorrelation unchanged:

\begin{equation}
T_{xx}=T_{yy}=\cos\theta,
\qquad
T_{xy}=-\sin\theta,
\qquad
T_{yx}=\sin\theta,
\qquad
T_{zz}=-1.
\label{eq:residual_phase_tensor_components}
\end{equation}

The transformation is local and unitary. Consequently,

\begin{equation}
\gamma=1,
\qquad
C=1,
\qquad
\tau=1,
\qquad
E_F=1
\label{eq:residual_phase_invariant_metrics}
\end{equation}

for every value of \( \theta \).

The fidelity with respect to the original Bell state is instead

\begin{equation}
F_{\Psi^+}
=
\frac{1+\cos\theta}{2}.
\label{eq:residual_phase_fidelity}
\end{equation}

Thus, a reduction in fixed-reference Bell-state fidelity does not by itself imply
entanglement degradation. It may result entirely from a coherent change in the
orientation of the state relative to the selected reference basis.

The residual phase model therefore establishes the first important reference case
of this thesis: coherent changes in fixed-basis measurements must be distinguished
from genuine decoherence and entanglement degradation.

%========================================================
% SECTION 2.2 — Statistical Phase Fluctuations
%========================================================

\section{Statistical Phase Fluctuations and Dephasing}
\label{sec:statistical_phase_fluctuations}

The fixed-phase description assumes that the relative phase is known and remains
constant over the relevant observation interval. In practice, birefringence and
optical-path conditions may vary between successive photon-pair realizations. If
these fluctuations are not resolved by the measurement system, the observed state
must be obtained by averaging over the corresponding phase realizations.

Let \(p_k\) denote the probability associated with the phase \( \theta_k \). The
ensemble density operator is

\begin{equation}
\rho_{AB}^{\mathrm{ens}}
=
\sum_k
p_k
|\psi(\theta_k)\rangle
\langle\psi(\theta_k)|,
\qquad
\sum_k p_k=1.
\label{eq:random_phase_ensemble}
\end{equation}

The phase statistics can be summarized through the complex coherence parameter

\begin{equation}
\mu
=
\left\langle e^{i\theta}\right\rangle
=
\sum_k p_k e^{i\theta_k}
=
\langle\cos\theta\rangle
+
i\langle\sin\theta\rangle.
\label{eq:coherence_parameter_mu}
\end{equation}

For a continuous probability density \(P(\theta)\),

\begin{equation}
\mu
=
\int d\theta\,
P(\theta)e^{i\theta}.
\label{eq:continuous_coherence_parameter_mu}
\end{equation}

It is useful to write

\begin{equation}
\mu
=
|\mu|e^{i\phi_\mu}.
\label{eq:coherence_parameter_polar_form}
\end{equation}

The magnitude \( |\mu| \) quantifies the coherence surviving the statistical
averaging, while the argument \( \phi_\mu \) represents the residual coherent
phase component of the state.

In the computational basis, the resulting state is

\begin{equation}
\rho_{AB}^{\mathrm{ens}}
=
\frac{1}{2}
\begin{pmatrix}
0 & 0 & 0 & 0\\
0 & 1 & \mu^* & 0\\
0 & \mu & 1 & 0\\
0 & 0 & 0 & 0
\end{pmatrix}.
\label{eq:random_phase_density_matrix}
\end{equation}

The diagonal populations remain unchanged, whereas the off-diagonal coherence
terms are multiplied by \( \mu \). The model therefore corresponds to a
random-unitary dephasing process together with a possible residual coherent phase
rotation when \( \arg(\mu)\neq0 \).

%--------------------------------------------------------
% Subsection 2.2.1 — Correlations and State Metrics
%--------------------------------------------------------

\subsection{Correlations and State Metrics}
\label{subsec:statistical_phase_metrics}

The correlation tensor becomes

\begin{equation}
T_\mu
=
\begin{pmatrix}
\mathrm{Re}(\mu) & -\mathrm{Im}(\mu) & 0\\
\mathrm{Im}(\mu) & \mathrm{Re}(\mu) & 0\\
0 & 0 & -1
\end{pmatrix}.
\label{eq:random_phase_correlation_tensor}
\end{equation}

The magnitude \( |\mu| \) therefore determines the surviving transverse
correlation amplitude, whereas the argument of \( \mu \) determines its remaining
orientation.

For this family of states,

\begin{equation}
\gamma
=
\frac{1+|\mu|^2}{2},
\label{eq:random_phase_purity}
\end{equation}

\begin{equation}
F_{\Psi^+}
=
\frac{1+\mathrm{Re}(\mu)}{2},
\label{eq:random_phase_fidelity}
\end{equation}

and

\begin{equation}
C
=
|\mu|,
\qquad
\tau
=
|\mu|^2.
\label{eq:random_phase_concurrence_tangle}
\end{equation}

The corresponding entanglement of formation is

\begin{equation}
E_F
=
h
\left(
\frac{
1+\sqrt{1-|\mu|^2}
}{2}
\right),
\label{eq:random_phase_eof}
\end{equation}

where \(h(x)\) is the binary entropy introduced in
Chapter~\ref{ch:01_mathematical_and_physical_framework}.

The limiting cases provide a direct physical interpretation. For a deterministic
phase,
\(
|\mu|=1,
\)
and the state remains pure and maximally entangled. 

For a phase uniformly distributed over \( [0,2\pi] \),
\(
\mu=0,
\)
and the transverse coherences vanish completely. The resulting state is

\begin{equation}
\rho_{\mu=0}
=
\frac{1}{2}
\left(
|01\rangle\langle01|
+
|10\rangle\langle10|
\right),
\label{eq:fully_dephased_bell_state}
\end{equation}

which retains perfect classical anticorrelation in the computational basis but is
no longer entangled.

Intermediate phase distributions produce partially coherent mixed states.

This model represents decoherence at the reduced polarization level, but it does
not yet specify the microscopic physical mechanism responsible for the unresolved
phase distribution. A frequency-resolved mechanism producing related losses of
polarization coherence is introduced later in this chapter.

%========================================================
% SECTION 2.3 — Depolarizing Reference Models
%========================================================

\section{Depolarizing Reference Models}
\label{sec:depolarizing_reference_models}

Statistical phase averaging suppresses specific coherence terms while preserving
the longitudinal population correlation. A complementary analytical reference is
provided by depolarizing channels, which describe effective loss of polarization
information directly at the density-operator level.

The depolarizing models introduced here are phenomenological. They do not
represent a microscopic model of PMD. Their purpose is to provide analytically
controlled descriptions of polarization degradation before the physical origin of
decoherence is treated explicitly through frequency-dependent propagation.

Two different classes of effective models are considered. Isotropic depolarization
contracts the polarization state uniformly along all Bloch-sphere directions,
whereas anisotropic depolarization preserves a preferred polarization axis and
suppresses only components orthogonal to it. The latter behavior is particularly
relevant as an effective reference for polarization-maintaining or PM-like fiber
configurations, where a physically distinguished birefringence axis is present.

%--------------------------------------------------------
% Subsection 2.3.1 — Global Depolarizing Baseline
%--------------------------------------------------------

\subsection{Global Depolarizing Baseline}
\label{subsec:global_depolarizing_baseline}

For a two-qubit density operator, the global depolarizing channel is defined as

\begin{equation}
\mathcal{D}^{4}_p(\rho)
=
(1-p)\rho
+
p\frac{I_4}{4},
\qquad
0\leq p\leq1.
\label{eq:global_two_qubit_depolarizing_channel}
\end{equation}

For a Bell-state input,

\begin{equation}
\rho(p)
=
(1-p)
|\Psi^+\rangle\langle\Psi^+|
+
p\frac{I_4}{4}.
\label{eq:globally_depolarized_bell_state}
\end{equation}

The corresponding purity is

\begin{equation}
\gamma
=
1-\frac{3}{2}p+\frac{3}{4}p^2,
\label{eq:global_depolarizing_purity}
\end{equation}

the Bell-state fidelity is

\begin{equation}
F_{\Psi^+}
=
1-\frac{3}{4}p,
\label{eq:global_depolarizing_fidelity}
\end{equation}

and the concurrence is

\begin{equation}
C
=
\max
\left(
0,\,
1-\frac{3}{2}p
\right).
\label{eq:global_depolarizing_concurrence}
\end{equation}

All bipartite Pauli correlations are contracted according to

\begin{equation}
T_{\mathrm{out}}
=
(1-p)T_{\mathrm{in}}.
\label{eq:global_depolarizing_tensor}
\end{equation}

The global channel is useful as a validation benchmark, but it is not directly
aligned with the geometry of the physical system considered in this thesis, where
the two photons propagate through separate fibers.

%--------------------------------------------------------
% Subsection 2.3.2 — Local Isotropic Depolarization
%--------------------------------------------------------

\subsection{Local Two-Arm Depolarization}
\label{subsec:local_two_arm_depolarization}

A description more closely aligned with the two-arm geometry assigns an
independent depolarizing channel to each propagation arm.

The parameter convention adopted throughout this work is defined directly by the
Bloch-vector contraction. A single-qubit isotropic depolarizing channel is written as

\begin{equation}
\mathcal{D}_{p}(\rho)
=
(1-p)\rho
+
p\frac{I_2}{2},
\label{eq:single_qubit_depolarizing_density_map}
\end{equation}

so that \(p=0\) corresponds to no depolarization and \(p=1\) to complete
depolarization.

Its linear extension to a generic operator \(X\) is

\begin{equation}
\mathcal{D}_{p}(X)
=
(1-p)X
+
p\frac{I_2}{2}\mathrm{Tr}(X).
\label{eq:single_qubit_depolarizing_operator_map}
\end{equation}

For a single-qubit state written in Bloch form,

\begin{equation}
\rho
=
\frac{1}{2}
\left(
I+\mathbf{r}\cdot\boldsymbol{\sigma}
\right),
\end{equation}

the channel gives

\begin{equation}
\mathcal{D}_p(\rho)
=
\frac{1}{2}
\left[
I+
(1-p)\mathbf{r}\cdot\boldsymbol{\sigma}
\right].
\label{eq:local_depolarizing_bloch_map}
\end{equation}

The Bloch vector therefore undergoes the isotropic contraction

\begin{equation}
\mathbf{r}
\longmapsto
(1-p)\mathbf{r}.
\label{eq:bloch_contraction_local_depolarization}
\end{equation}

For two independent fiber arms,

\begin{equation}
\rho_{AB}^{\mathrm{out}}
=
\left(
\mathcal{D}_{p_A}
\otimes
\mathcal{D}_{p_B}
\right)
\left(
\rho_{AB}^{\mathrm{in}}
\right).
\label{eq:two_arm_local_depolarizing_channel}
\end{equation}

Although the effective noise acts locally, the channel acts on the complete
bipartite density operator. The entanglement and two-photon correlations are
properties of the joint state and cannot be reconstructed from the two reduced
states alone.

%--------------------------------------------------------
% Subsection 2.3.3 — Anisotropic PM-Like Depolarization
%--------------------------------------------------------

\subsection{Anisotropic Depolarization for a PM-Like Fiber}
\label{subsec:anisotropic_pm_depolarization}

The isotropic model assumes that polarization information is degraded equally
along all directions of the Bloch sphere. This assumption is not appropriate when
the propagation medium possesses a physically preferred polarization axis.

A polarization-maintaining or PM-like fiber provides a natural example. Its two
principal polarization modes define a distinguished axis in polarization space.
When propagation couples polarization to an unresolved degree of freedom, such as
frequency or arrival time, coherence between the corresponding principal
polarization modes may be reduced while polarization information along the
principal axis is preserved.

This behavior is described by an anisotropic dephasing channel. In the standard
form commonly introduced for a preferred \(Z\) axis, the corresponding Pauli
channel preserves the \(z\)-component of the Bloch vector while contracting the
transverse \(x\)- and \(y\)-components. An equivalent channel around any other
axis is obtained through an appropriate unitary change of basis.

For the convention adopted in this work, the preferred PM-like axis is chosen to
be the \(X\) axis. The corresponding effective channel is defined as

\begin{equation}
\mathcal{A}^{(x)}_{p}(\rho)
=
\left(
1-\frac{p}{2}
\right)\rho
+
\frac{p}{2}
\sigma_x\rho\sigma_x,
\qquad
0\leq p\leq1.
\label{eq:anisotropic_x_depolarizing_channel}
\end{equation}

This parametrization is chosen so that \(p\) has the same operational convention
used for the isotropic channel: \(p=0\) corresponds to no degradation, while
\(p=1\) corresponds to complete suppression of the Bloch components orthogonal
to the preferred axis.

The map in Eq.~\eqref{eq:anisotropic_x_depolarizing_channel} is a probabilistic
Pauli-\(X\) channel. Its action on the Pauli basis is

\begin{equation}
\mathcal{A}^{(x)}_{p}(I)
=
I,
\end{equation}

\begin{equation}
\mathcal{A}^{(x)}_{p}(\sigma_x)
=
\sigma_x,
\label{eq:anisotropic_x_pauli_preserved}
\end{equation}

and

\begin{equation}
\mathcal{A}^{(x)}_{p}(\sigma_y)
=
(1-p)\sigma_y,
\qquad
\mathcal{A}^{(x)}_{p}(\sigma_z)
=
(1-p)\sigma_z.
\label{eq:anisotropic_x_pauli_contracted}
\end{equation}

Therefore, for a Bloch vector

\begin{equation}
\mathbf{r}
=
\begin{pmatrix}
r_x\\
r_y\\
r_z
\end{pmatrix},
\end{equation}

the channel produces

\begin{equation}
\mathbf{r}
\longmapsto
\begin{pmatrix}
r_x\\
(1-p)r_y\\
(1-p)r_z
\end{pmatrix}.
\label{eq:anisotropic_x_bloch_contraction}
\end{equation}

The Bloch sphere is thus contracted into an ellipsoid whose major axis coincides
with the preferred \(X\) direction. States lying on this axis are unaffected,
whereas components orthogonal to it progressively lose coherence.

This behavior differs fundamentally from isotropic depolarization. In the
isotropic case,

\begin{equation}
(r_x,r_y,r_z)
\longmapsto
(1-p)(r_x,r_y,r_z),
\end{equation}

whereas the anisotropic PM-like model gives

\begin{equation}
(r_x,r_y,r_z)
\longmapsto
\left(
r_x,\,
(1-p)r_y,\,
(1-p)r_z
\right).
\end{equation}

The choice of the \(X\) axis is a convention rather than a restriction of the
model. For a generic unit vector \(\mathbf{n}\), the corresponding preferred-axis
Pauli operator is

\begin{equation}
\sigma_{\mathbf{n}}
=
\mathbf{n}\cdot\boldsymbol{\sigma},
\end{equation}

and the anisotropic channel can be written as

\begin{equation}
\mathcal{A}^{(\mathbf{n})}_{p}(\rho)
=
\left(
1-\frac{p}{2}
\right)\rho
+
\frac{p}{2}
\sigma_{\mathbf{n}}
\rho
\sigma_{\mathbf{n}}.
\label{eq:anisotropic_general_axis_channel}
\end{equation}

In Bloch-vector form, decomposing

\begin{equation}
\mathbf{r}
=
\mathbf{r}_{\parallel}
+
\mathbf{r}_{\perp},
\end{equation}

with

\begin{equation}
\mathbf{r}_{\parallel}
=
(\mathbf{r}\cdot\mathbf{n})\mathbf{n},
\qquad
\mathbf{r}_{\perp}
=
\mathbf{r}
-
\mathbf{r}_{\parallel},
\end{equation}

the channel becomes

\begin{equation}
\mathbf{r}
\longmapsto
\mathbf{r}_{\parallel}
+
(1-p)\mathbf{r}_{\perp}.
\label{eq:anisotropic_general_bloch_map}
\end{equation}

This expression makes the physical interpretation explicit: polarization
information parallel to the preferred axis survives, whereas coherence in the
orthogonal plane is progressively suppressed.

For two independent PM-like fiber arms with the same preferred \(X\) axis,

\begin{equation}
\rho_{AB}^{\mathrm{out}}
=
\left(
\mathcal{A}^{(x)}_{p_A}
\otimes
\mathcal{A}^{(x)}_{p_B}
\right)
\left(
\rho_{AB}^{\mathrm{in}}
\right).
\label{eq:two_arm_anisotropic_x_channel}
\end{equation}

For the reference Bell state \( |\Psi^+\rangle \), whose correlation tensor is

\begin{equation}
T_{\Psi^+}
=
\mathrm{diag}
\left(
1,\,
1,\,
-1
\right),
\label{eq:psi_plus_correlation_tensor_pm}
\end{equation}

the anisotropic two-arm channel gives

\begin{equation}
T_{\mathrm{out}}
=
\mathrm{diag}
\left(
1,\,
\eta_{\perp},\,
-\eta_{\perp}
\right),
\label{eq:anisotropic_pm_tensor_psiplus}
\end{equation}

with

\begin{equation}
\eta_{\perp}
=
(1-p_A)(1-p_B).
\label{eq:anisotropic_transverse_survival_factor}
\end{equation}

Consequently,

\begin{equation}
\langle
\sigma_x\otimes\sigma_x
\rangle_{\mathrm{out}}
=
1,
\end{equation}

while

\begin{equation}
\langle
\sigma_y\otimes\sigma_y
\rangle_{\mathrm{out}}
=
\eta_{\perp},
\qquad
\langle
\sigma_z\otimes\sigma_z
\rangle_{\mathrm{out}}
=
-\eta_{\perp}.
\label{eq:anisotropic_psiplus_correlations}
\end{equation}

The preferred-axis correlation therefore survives completely, while correlations
in the orthogonal directions are attenuated.

This effective channel should not be interpreted as the microscopic propagation
law of a PM fiber. Its role is instead to provide a polarization-only reference
model for the anisotropic decoherence that can emerge when a preferred
birefringence axis is present and unresolved spectral or temporal degrees of
freedom are traced out. The corresponding microscopic description will be
developed later through frequency-dependent birefringent propagation and PMD.

%--------------------------------------------------------
% Subsection 2.3.4 — Correlation-Survival Factor
%--------------------------------------------------------

\subsection{Correlation-Survival Factor}
\label{subsec:correlation_survival_factor}

For the isotropic local depolarizing channel, the single-qubit action on the Pauli
basis is

\begin{equation}
\mathcal{D}_p(I)=I,
\qquad
\mathcal{D}_p(\sigma_i)
=
(1-p)\sigma_i,
\qquad
i\in\{x,y,z\}.
\label{eq:depolarizing_pauli_action}
\end{equation}

Consequently, each two-body Pauli correlation is multiplied by the common factor

\begin{equation}
\eta
=
(1-p_A)(1-p_B).
\label{eq:correlation_survival_factor}
\end{equation}

The correlation tensor therefore satisfies

\begin{equation}
T_{\mathrm{out}}
=
\eta T_{\mathrm{in}},
\label{eq:local_depolarizing_tensor_contraction}
\end{equation}

and, for arbitrary local measurement axes,

\begin{equation}
E_{\mathrm{out}}(\mathbf{a},\mathbf{b})
=
\eta
E_{\mathrm{in}}(\mathbf{a},\mathbf{b}).
\label{eq:local_depolarizing_correlation_contraction}
\end{equation}

The factor \( \eta \) will be used throughout the isotropic effective fiber-channel
models as the survival factor of bipartite polarization correlations.

For anisotropic depolarization, however, no single scalar survival factor applies
to all correlation components. Instead, the contraction is axis dependent. For
the \(X\)-oriented PM-like model introduced in
Section~\ref{subsec:anisotropic_pm_depolarization}, the \(x\)-component is
preserved while the transverse \(y\)- and \(z\)-components acquire the factor
\(\eta_{\perp}\). The isotropic and anisotropic channels must therefore be kept
distinct in the subsequent analytical and numerical treatments.

%========================================================
% SECTION 2.4 — Composite Effective Fiber Channel
%========================================================

\section{Composite Effective Fiber Channel}
\label{sec:composite_effective_fiber_channel}

The fixed relative-phase model and the local depolarizing model represent two
different classes of propagation effects. The former is coherent and rotates the
correlation structure, whereas the latter is non-unitary at the polarization level
and suppresses its magnitude.

They are combined through the effective channel

\begin{equation}
\rho_{AB}^{\mathrm{out}}
=
\left(
\mathcal{D}_{p_A}
\otimes
\mathcal{D}_{p_B}
\right)
\left[
U_\theta
\rho_{AB}^{\mathrm{in}}
U_\theta^\dagger
\right].
\label{eq:composite_fiber_channel}
\end{equation}

For isotropic depolarization, the local depolarizing channel is covariant under
single-qubit unitary transformations. Consequently, for the phase transformation
considered here,

\begin{equation}
\left(
\mathcal{D}_{p_A}\otimes\mathcal{D}_{p_B}
\right)
\left[
U_\theta\rho U_\theta^\dagger
\right]
=
U_\theta
\left[
\left(
\mathcal{D}_{p_A}\otimes\mathcal{D}_{p_B}
\right)(\rho)
\right]
U_\theta^\dagger.
\label{eq:phase_depolarization_commutation}
\end{equation}

Thus, within this isotropic model, the order of phase evolution and
depolarization does not affect the final state. This property is specific to the
adopted channel and does not hold for arbitrary anisotropic noise mechanisms.

%--------------------------------------------------------
% Subsection 2.4.1 — Fixed-Phase Composite Model
%--------------------------------------------------------

\subsection{Fixed-Phase Composite Model}
\label{subsec:fixed_phase_composite_model}

For the maximally entangled input family considered here, the reduced states of
both subsystems remain \(I_2/2\). The action of the two local depolarizing channels
can therefore be written compactly as

\begin{equation}
\rho_{AB}^{\mathrm{out}}
=
\eta
|\psi(\theta)\rangle\langle\psi(\theta)|
+
(1-\eta)\frac{I_4}{4}.
\label{eq:fixed_composite_density_matrix_compact}
\end{equation}

The corresponding correlation tensor is

\begin{equation}
T_{\mathrm{out}}
=
\eta
\begin{pmatrix}
\cos\theta & -\sin\theta & 0\\
\sin\theta & \cos\theta & 0\\
0 & 0 & -1
\end{pmatrix}.
\label{eq:fixed_composite_tensor}
\end{equation}

The parameter \( \theta \) therefore controls the orientation of the transverse
correlations, whereas \( \eta \) controls their magnitude.

The corresponding purity is

\begin{equation}
\gamma
=
\frac{1+3\eta^2}{4},
\label{eq:fixed_composite_purity}
\end{equation}

while the Bell-state fidelity is

\begin{equation}
F_{\Psi^+}
=
\frac{1-\eta}{4}
+
\eta
\frac{1+\cos\theta}{2}.
\label{eq:fixed_composite_fidelity}
\end{equation}

The concurrence is

\begin{equation}
C
=
\max
\left(
0,\,
\frac{3\eta-1}{2}
\right).
\label{eq:fixed_composite_concurrence}
\end{equation}

%--------------------------------------------------------
% Subsection 2.4.2 — Statistical Composite Model
%--------------------------------------------------------

\subsection{Statistical Composite Model}
\label{subsec:statistical_composite_model}

The composite description can be extended by allowing the relative phase to
fluctuate between realizations. The ensemble-averaged state is

\begin{equation}
\rho_{AB}^{\mathrm{ens}}
=
\sum_k p_k
\left(
\mathcal{D}_{p_A}
\otimes
\mathcal{D}_{p_B}
\right)
\left[
U_{\theta_k}
\rho_{AB}^{\mathrm{in}}
U_{\theta_k}^\dagger
\right].
\label{eq:statistical_composite_channel}
\end{equation}

Using the covariance property in
Eq.~\eqref{eq:phase_depolarization_commutation},

\begin{equation}
\rho_{AB}^{\mathrm{ens}}
=
\left(
\mathcal{D}_{p_A}
\otimes
\mathcal{D}_{p_B}
\right)
\left[
\sum_k
p_k
U_{\theta_k}
\rho_{AB}^{\mathrm{in}}
U_{\theta_k}^\dagger
\right].
\label{eq:statistical_composite_reordered}
\end{equation}

For the random-phase family of Section~\ref{sec:statistical_phase_fluctuations},
this reduces to

\begin{equation}
\rho_{AB}^{\mathrm{ens}}
=
\eta
\rho_{AB}^{\mu}
+
(1-\eta)
\frac{I_4}{4},
\label{eq:statistical_composite_density_compact}
\end{equation}

where \( \rho_{AB}^{\mu} \) denotes the state in
Eq.~\eqref{eq:random_phase_density_matrix}.

The output correlation tensor is

\begin{equation}
T_{\mathrm{out}}^{\mathrm{ens}}
=
\eta
\begin{pmatrix}
\mathrm{Re}(\mu) & -\mathrm{Im}(\mu) & 0\\
\mathrm{Im}(\mu) & \mathrm{Re}(\mu) & 0\\
0 & 0 & -1
\end{pmatrix}.
\label{eq:statistical_composite_tensor}
\end{equation}

The two mechanisms therefore leave distinguishable signatures. Phase averaging
suppresses the transverse coherence through \( \mu \), whereas local isotropic
depolarization multiplies all bipartite correlations by \( \eta \).

The purity becomes

\begin{equation}
\gamma
=
\frac{
1+\eta^2
\left(
1+2|\mu|^2
\right)
}{4},
\label{eq:statistical_composite_purity}
\end{equation}

the fidelity is

\begin{equation}
F_{\Psi^+}
=
\frac{1-\eta}{4}
+
\eta
\frac{
1+\mathrm{Re}(\mu)
}{2},
\label{eq:statistical_composite_fidelity}
\end{equation}

and the concurrence is

\begin{equation}
C
=
\max
\left(
0,\,
\eta|\mu|
-
\frac{1-\eta}{2}
\right).
\label{eq:statistical_composite_concurrence}
\end{equation}

At the level of transverse polarization correlations, the progression of the
effective models can therefore be summarized as

\begin{equation}
\boxed{
e^{i\theta}
\quad\longrightarrow\quad
\mu=\left\langle e^{i\theta}\right\rangle
\quad\longrightarrow\quad
\eta\mu
}
\label{eq:effective_model_progression}
\end{equation}

corresponding respectively to coherent phase evolution, statistical dephasing,
and the additional suppression produced by isotropic depolarization.
The effective model is analytically transparent and useful for distinguishing
coherent phase evolution, statistical dephasing, and isotropic polarization
degradation. Its main limitation is that the parameters
\(\mu\), \(p_A\), and \(p_B\) describe phenomenological behavior rather than being derived
directly from the spectral and birefringent properties of the physical fibers.

%========================================================
% SECTION 2.5 — Frequency-Dependent Polarization Propagation
%========================================================

\section{Frequency-Dependent Polarization Propagation}
\label{sec:frequency_dependent_propagation}

The effective models introduced in the previous sections treat each photon
exclusively as a polarization qubit. This description is sufficient when the
polarization transformation can be considered approximately constant over the
photon bandwidth. For finite-bandwidth photons propagating through a birefringent
fiber, however, different spectral components may experience different
polarization transformations.

The local polarization operator must therefore be generalized to

\begin{equation}
U_K
\longrightarrow
U_K(\omega_K),
\qquad
K\in\{A,B\}.
\label{eq:frequency_dependent_local_unitary}
\end{equation}

For a resolved pair of frequencies, the corresponding two-arm transformation is

\begin{equation}
U_{AB}(\omega_A,\omega_B)
=
U_A(\omega_A)
\otimes
U_B(\omega_B).
\label{eq:frequency_dependent_two_arm_unitary}
\end{equation}

The complete polarization--frequency evolution of an ideal lossless fiber remains
unitary. Nevertheless, when \(U_K(\omega_K)\) varies over the occupied bandwidth,
different spectral components become associated with different polarization
transformations. Polarization and frequency can therefore become correlated during
propagation.

If the spectral degree of freedom is not resolved by the measurement system, only
the reduced polarization state is observed. The resulting polarization evolution
may then appear non-unitary even though the complete optical evolution remains
unitary.

This provides the physical connection with the effective models introduced
previously: polarization decoherence can emerge from frequency-dependent coherent
propagation followed by the neglect of an unresolved spectral degree of freedom.


%========================================================
% SECTION 2.6 — Spectral Description of the Photon Pair
%========================================================

\section{Spectral Description of the Photon Pair}
\label{sec:spectral_description_spdc}

To describe frequency-dependent propagation, the Hilbert space associated with
each photon must include both spectral and polarization degrees of freedom,

\begin{equation}
\mathcal{H}_K
=
\mathcal{H}_{\omega_K}
\otimes
\mathcal{H}_{\mathrm{pol},K},
\qquad
K\in\{A,B\}.
\label{eq:single_photon_frequency_polarization_space}
\end{equation}

The complete two-photon space can consequently be written as

\begin{equation}
\mathcal{H}_{\mathrm{tot}}
=
\left(
\mathcal{H}_{\omega_A}
\otimes
\mathcal{H}_{\omega_B}
\right)
\otimes
\left(
\mathcal{H}_{\mathrm{pol},A}
\otimes
\mathcal{H}_{\mathrm{pol},B}
\right).
\label{eq:complete_frequency_polarization_hilbert_space}
\end{equation}

The finite-bandwidth spectral state of the photon pair is represented as

\begin{equation}
|F\rangle_{\omega_A\omega_B}
=
\iint
d\omega_A\,d\omega_B\,
F(\omega_A,\omega_B)
|\omega_A,\omega_B\rangle,
\label{eq:spdc_two_photon_spectral_state}
\end{equation}

where \(F(\omega_A,\omega_B)\) is the joint spectral amplitude (JSA). It is
normalized according to

\begin{equation}
\iint
d\omega_A\,d\omega_B\,
|F(\omega_A,\omega_B)|^2
=
1.
\label{eq:joint_spectral_amplitude_normalization}
\end{equation}

The corresponding joint spectral intensity (JSI) is

\begin{equation}
J(\omega_A,\omega_B)
=
|F(\omega_A,\omega_B)|^2,
\label{eq:joint_spectral_intensity}
\end{equation}

and represents the joint probability density of the two photon frequencies.

For an SPDC source, the joint spectral structure is determined by the pump
spectrum and the phase-matching properties of the nonlinear medium. The JSA may
be written schematically as

\begin{equation}
F(\omega_A,\omega_B)
=
\mathcal{N}\,
\alpha_p(\omega_A+\omega_B)
\Phi(\omega_A,\omega_B),
\label{eq:spdc_joint_spectral_amplitude}
\end{equation}

where \(\alpha_p\) is the pump spectral envelope, \(\Phi\) is the phase-matching
function, and \(\mathcal{N}\) is a normalization factor
~\cite{grice_walmsley_spdc_spectral_1997}.

The resulting frequencies may be statistically correlated. In particular, a
narrow pump spectrum constrains the sum frequency
\(\omega_A+\omega_B\), producing the frequency anticorrelation commonly observed
in SPDC photon pairs.


%--------------------------------------------------------
% Subsection 2.6.1 — Input Polarization--Frequency State
%--------------------------------------------------------

\subsection{Input Polarization--Frequency State}
\label{subsec:input_polarization_frequency_state}

The model developed in this work assumes that the spectral and polarization
degrees of freedom are initially separable. In particular, the polarization state
is assumed to be the same for every generated frequency pair.

For the reference Bell-state input,

\begin{equation}
|\Psi_{\mathrm{tot}}^{\mathrm{in}}\rangle
=
|F\rangle_{\omega_A\omega_B}
\otimes
|\Psi^+\rangle_{\mathrm{pol}}.
\label{eq:factorized_spectral_polarization_input_ket}
\end{equation}

More generally,

\begin{equation}
\rho_{\mathrm{tot}}^{\mathrm{in}}
=
|F\rangle\langle F|
\otimes
\rho_{AB}^{\mathrm{pol,in}}.
\label{eq:factorized_spectral_polarization_input}
\end{equation}

This factorization is an explicit modeling assumption. It neglects initial
polarization--frequency distinguishability at the source, so that any coupling
between these degrees of freedom appearing after transmission is attributed to
the frequency dependence of the fiber propagation
~\cite{poon_law_stochastic_pmd_2008,
riccardi_polarization_entanglement_distribution_2021}.


%--------------------------------------------------------
% Subsection 2.6.2 — Frequency-Dependent Two-Photon Evolution
%--------------------------------------------------------

\subsection{Frequency-Dependent Two-Photon Evolution}
\label{subsec:frequency_dependent_two_photon_evolution}

Each spectral component of the input state undergoes the corresponding local
polarization transformations,

\begin{equation}
U_{AB}(\omega_A,\omega_B)
=
U_A(\omega_A)
\otimes
U_B(\omega_B).
\label{eq:frequency_dependent_two_photon_operator}
\end{equation}

After propagation, the complete polarization--frequency state remains unitary.
If frequency is not resolved, the observable polarization state is obtained by
tracing over the spectral degrees of freedom,

\begin{equation}
\rho_{AB}^{\mathrm{pol,out}}
=
\mathrm{Tr}_{\omega_A,\omega_B}
\left[
\rho_{\mathrm{tot}}^{\mathrm{out}}
\right].
\label{eq:two_photon_spectral_partial_trace}
\end{equation}

Using the factorized input state of
Eq.~\eqref{eq:factorized_spectral_polarization_input}, the reduced polarization
state becomes

\begin{equation}
\boxed{
\rho_{AB}^{\mathrm{pol,out}}
=
\iint
d\omega_A\,d\omega_B\,
J(\omega_A,\omega_B)
U_{AB}(\omega_A,\omega_B)
\rho_{AB}^{\mathrm{pol,in}}
U_{AB}^{\dagger}(\omega_A,\omega_B)
}
\label{eq:continuous_reduced_polarization_state}
\end{equation}

~\cite{poon_law_stochastic_pmd_2008,
antonelli_shtaif_brodsky_pmd_2011}.

Equation~\eqref{eq:continuous_reduced_polarization_state} provides the central
connection between frequency-dependent fiber propagation and the effective
polarization state. The unresolved spectrum produces a continuous average over
the polarization transformations experienced by the different frequency pairs.

The resulting reduced evolution is not in general equivalent to the isotropic
depolarizing channel introduced in
Section~\ref{sec:depolarizing_reference_models}. Its form depends on the
frequency dependence of the fiber transformation. For fixed birefringence axes,
for example, the reduced evolution has the form of dephasing in the corresponding
eigenpolarization basis.


%========================================================
% SECTION 2.7 — Birefringence and Polarization-Mode Dispersion
%========================================================

\section{Birefringence and Polarization-Mode Dispersion}
\label{sec:birefringence_and_pmd}

The frequency dependence of the polarization transformation originates from fiber
birefringence. A locally uniform birefringent segment supports two orthogonal
polarization eigenmodes, conventionally referred to as slow and fast, with
propagation constants \(\beta_s(\omega)\) and \(\beta_f(\omega)\)
~\cite{gordon_kogelnik_pmd_2000,
wai_menyuk_random_birefringence_1996}.

Their differential propagation constant is

\begin{equation}
\Delta\beta(\omega)
=
\beta_s(\omega)-\beta_f(\omega).
\label{eq:differential_propagation_constant}
\end{equation}

Using the corresponding effective refractive indices,

\begin{equation}
\beta_{s,f}(\omega)
=
\frac{\omega}{c}
n_{s,f}(\omega),
\end{equation}

one obtains

\begin{equation}
\Delta\beta(\omega)
=
\frac{\omega}{c}
\Delta n_{\mathrm{eff}}(\omega),
\qquad
\Delta n_{\mathrm{eff}}
=
n_s-n_f.
\label{eq:differential_beta_refractive_index}
\end{equation}

For a uniform segment of length \(L_j\), the accumulated differential phase is

\begin{equation}
\Delta\phi_j(\omega)
=
L_j\Delta\beta_j(\omega).
\label{eq:differential_phase_beta_relation}
\end{equation}

Since this phase depends on frequency, the two eigenmodes experience different
group delays. The differential group delay (DGD) is determined by

\begin{equation}
\Delta\tau_j
=
\left.
\frac{d\Delta\phi_j}{d\omega}
\right|_{\omega_0}
=
L_j
\left.
\frac{d\Delta\beta_j}{d\omega}
\right|_{\omega_0}.
\label{eq:dgd_phase_derivative}
\end{equation}

Its magnitude \( |\Delta\tau_j| \) quantifies the temporal separation accumulated
between the two local polarization eigenmodes. The physical relation can therefore
be summarized as

\begin{equation}
\boxed{
\Delta n_{\mathrm{eff}}
\longrightarrow
\Delta\beta
\longrightarrow
\Delta\phi(\omega)
\longrightarrow
\Delta\tau
}
\label{eq:birefringence_to_dgd_chain}
\end{equation}


%--------------------------------------------------------
% Subsection 2.7.1 — Single Birefringent Segment
%--------------------------------------------------------

\subsection{Single Birefringent Segment}
\label{subsec:single_birefringent_segment}

Consider a locally uniform, linear, and lossless birefringent fiber segment. Its
polarization evolution can be represented by the frequency-dependent unitary
Jones operator

\begin{equation}
U_j(\omega)
=
e^{i\bar{\phi}_j(\omega)}
\exp
\left[
\frac{i}{2}
\Delta\phi_j(\omega)
\hat{\mathbf{n}}_j
\cdot
\boldsymbol{\sigma}
\right],
\label{eq:single_birefringent_segment_operator}
\end{equation}

where \(\bar{\phi}_j(\omega)\) is a polarization-independent phase,
\(\Delta\phi_j(\omega)\) is the differential phase, and
\(\hat{\mathbf{n}}_j\) identifies the birefringence axis on the Bloch sphere.

Around the carrier frequency \(\omega_0\), the differential phase can be expanded
to first order as

\begin{equation}
\Delta\phi_j(\omega)
\simeq
\Delta\phi_j(\omega_0)
+
\Delta\tau_j
(\omega-\omega_0).
\label{eq:first_order_differential_phase}
\end{equation}

The first term describes the coherent birefringent phase at the carrier
frequency, whereas the second describes its first-order frequency dependence.
The latter is responsible for first-order PMD.

At this level, \(\hat{\mathbf n}_j\) defines the two orthogonal principal
polarization directions of the segment, while \( |\Delta\tau_j| \) gives their
differential group delay.


%--------------------------------------------------------
% Subsection 2.7.2 — Fixed-Frequency Limit and Residual Phase Model
%--------------------------------------------------------

\subsection{Fixed-Frequency Limit and Connection to the Residual Phase Model}
\label{subsec:fixed_frequency_birefringence_residual_phase}

At a fixed frequency, the transformation in
Eq.~\eqref{eq:single_birefringent_segment_operator} is a coherent unitary
polarization transformation. In the local birefringence eigenbasis it can be
written, apart from a global phase, as

\begin{equation}
U_j^{\mathrm{eig}}(\omega_0)
\sim
\begin{pmatrix}
e^{-i\Delta\phi_j/2} & 0\\
0 & e^{+i\Delta\phi_j/2}
\end{pmatrix}.
\label{eq:fixed_frequency_eigenbasis_operator}
\end{equation}

If the relevant eigenbasis is aligned with the computational polarization basis,
factoring out the phase of the first component gives

\begin{equation}
U_j^{\mathrm{eig}}(\omega_0)
\sim
\begin{pmatrix}
1 & 0\\
0 & e^{i\Delta\phi_j}
\end{pmatrix}.
\label{eq:fixed_frequency_reduced_phase_form}
\end{equation}

This is the same operator form introduced in the residual phase model of
Section~\ref{sec:reduced_phase_model}.

For a complete fiber, the local birefringence axes need not coincide with the
laboratory \(H/V\) basis and may vary along the propagation direction. The
fixed-frequency input--output transformation is then obtained from the product of
the local segment transformations. After polarization compensation, the reduced
model introduced earlier represents the remaining coherent transformation as

\begin{equation}
U_{\mathrm{fib}}^{\mathrm{comp}}(\omega_0)
\sim
\begin{pmatrix}
1 & 0\\
0 & e^{i\theta}
\end{pmatrix},
\label{eq:compensated_complete_fiber_reduced_operator}
\end{equation}

where \(\theta\) is the residual coherent phase not removed by the compensation
procedure.

The residual phase model is therefore an effective fixed-frequency description of
the compensated fiber transformation. At fixed frequency this evolution remains
unitary and cannot by itself reduce polarization purity or entanglement.
Decoherence emerges when the frequency dependence of the transformation becomes
relevant over the finite photon bandwidth and the unresolved spectral degree of
freedom is traced out.

%--------------------------------------------------------
% Subsection 2.7.2 — First-Order PMD of a Single Segment
%--------------------------------------------------------

\subsection{First-Order PMD of a Single Segment}
\label{subsec:first_order_single_segment_pmd}

It is convenient to introduce the angular-frequency offset from the carrier,

\begin{equation}
\Omega
=
\omega-\omega_0.
\label{eq:angular_frequency_offset}
\end{equation}

Using the first-order expansion of
Eq.~\eqref{eq:first_order_differential_phase}, the polarization-dependent part
of the segment transformation becomes

\begin{equation}
\mathcal{U}_j(\Omega)
\simeq
\exp
\left[
\frac{i}{2}
\left(
\Delta\phi_j(\omega_0)
+
\Omega\Delta\tau_j
\right)
\hat{\mathbf n}_j
\cdot
\boldsymbol{\sigma}
\right].
\label{eq:segment_first_order_frequency_operator}
\end{equation}

If the polarization transformation at the carrier frequency is compensated, the
remaining first-order frequency-dependent transformation is

\begin{equation}
\boxed{
\mathcal{U}_j^{\mathrm{res}}(\Omega)
=
\exp
\left[
\frac{i}{2}
\Omega\Delta\tau_j
\hat{\mathbf n}_j
\cdot
\boldsymbol{\sigma}
\right].
}
\label{eq:residual_single_segment_operator}
\end{equation}

Let \( |s_+\rangle \) and \( |s_-\rangle \) denote the eigenstates of the local
birefringence generator,

\begin{equation}
\left(
\hat{\mathbf n}_j
\cdot
\boldsymbol{\sigma}
\right)
|s_\pm\rangle
=
\pm |s_\pm\rangle.
\label{eq:segment_birefringence_generator_eigenstates}
\end{equation}

An arbitrary input polarization can be written as

\begin{equation}
|\chi\rangle
=
\alpha|s_+\rangle
+
\beta|s_-\rangle.
\label{eq:input_segment_eigenmode_superposition}
\end{equation}

For a finite-bandwidth photon with spectral amplitude \(f_\Omega(\Omega)\), the
propagated state is

\begin{equation}
|\Psi_{\mathrm{out}}\rangle
=
\int d\Omega\,
f_\Omega(\Omega)
|\omega_0+\Omega\rangle
\otimes
\left[
\alpha
e^{+i\Omega\Delta\tau_j/2}
|s_+\rangle
+
\beta
e^{-i\Omega\Delta\tau_j/2}
|s_-\rangle
\right].
\label{eq:single_segment_frequency_polarization_state}
\end{equation}

The two eigenpolarization components therefore acquire opposite
frequency-dependent phases. In the time domain, this corresponds to two
wave-packet components separated by the differential group delay
\( |\Delta\tau_j| \), as illustrated in
Fig.~\ref{fig:pmd_temporal_splitting}.

\begin{figure}[htbp]
    \centering

    \begin{tikzpicture}[
        >=Latex,
        every node/.style={font=\small}
    ]

        % Input axis
        \draw[->] (0,0) -- (2.7,0)
        node[below] {\(t\)};

        % Input pulse
        \draw[thick, domain=0.30:2.35, samples=70]
        plot (\x,{1.15*exp(-4.2*(\x-1.30)^2)});

        \node at (1.30,1.55)
        {input superposition};

        % Arrow into fiber
        \draw[->, thick] (2.65,0.55) -- (3.35,0.55);

        % Fiber segment
        \node[
            draw,
            rounded corners,
            minimum width=2.3cm,
            minimum height=1.5cm,
            align=center
        ] at (4.55,0.55)
        {birefringent\\segment};

        \node at (4.55,-0.55)
        {\(\hat{\mathbf n}_j,\ \Delta\tau_j\)};

        % Arrow out
        \draw[->, thick] (5.75,0.55) -- (6.45,0.55);

        % Output time axis
        \draw[->] (6.45,0) -- (10.15,0)
        node[below] {\(t\)};

        % Output + component
        \draw[thick, domain=6.55:8.50, samples=70]
        plot (\x,{0.95*exp(-5.0*(\x-7.45)^2)});

        % Output - component
        \draw[thick, dashed, domain=8.05:10.00, samples=70]
        plot (\x,{0.95*exp(-5.0*(\x-9.10)^2)});

        \node at (7.45,1.30)
        {\(|s_+\rangle\)};

        \node at (9.10,1.30)
        {\(|s_-\rangle\)};

        % DGD marker
        \draw[<->]
        (7.45,-0.50)
        --
        (9.10,-0.50)
        node[midway,below]
        {\(|\Delta\tau_j|\)};

    \end{tikzpicture}

    \caption{
        Time-domain interpretation of first-order PMD in a single birefringent
        segment. The two local eigenpolarization components are separated by the
        differential group delay \( |\Delta\tau_j| \).
    }
    \label{fig:pmd_temporal_splitting}
\end{figure}

If frequency is not resolved, tracing over the spectral degree of freedom gives

\begin{equation}
\rho_{\mathrm{pol}}^{\mathrm{out}}
=
\operatorname{Tr}_{\omega}
\left[
|\Psi_{\mathrm{out}}\rangle
\langle\Psi_{\mathrm{out}}|
\right].
\label{eq:single_segment_spectral_partial_trace}
\end{equation}

In the local eigenpolarization basis,

\begin{equation}
\rho_{\mathrm{pol}}^{\mathrm{out}}
=
\begin{pmatrix}
\rho_{++}^{\mathrm{in}}
&
\kappa_j\rho_{+-}^{\mathrm{in}}
\\
\kappa_j^*\rho_{-+}^{\mathrm{in}}
&
\rho_{--}^{\mathrm{in}}
\end{pmatrix},
\label{eq:single_segment_reduced_polarization_state}
\end{equation}

where

\begin{equation}
\boxed{
\kappa_j
=
\int d\Omega\,
|f_\Omega(\Omega)|^2
e^{i\Omega\Delta\tau_j}.
}
\label{eq:single_segment_coherence_factor}
\end{equation}

The magnitude \( |\kappa_j| \) therefore quantifies the polarization coherence
remaining after the unresolved spectral degree of freedom has been traced out
~\cite{poon_law_stochastic_pmd_2008}.

For a Gaussian spectral probability density,

\begin{equation}
|f_\Omega(\Omega)|^2
=
\frac{1}{\sqrt{2\pi}\sigma_\omega}
\exp
\left[
-\frac{\Omega^2}{2\sigma_\omega^2}
\right],
\label{eq:gaussian_spectral_probability_density}
\end{equation}

the coherence integral in
Eq.~\eqref{eq:single_segment_coherence_factor} is the characteristic function of
the Gaussian spectral distribution and gives

\begin{equation}
\boxed{
|\kappa_j|
=
\exp
\left[
-\frac{1}{2}
\sigma_\omega^2
\Delta\tau_j^2
\right].
}
\label{eq:gaussian_pmd_coherence_factor}
\end{equation}

This Gaussian decay of polarization coherence with increasing spectral width and
differential group delay is the analytical single-segment limit of the
frequency-dependent PMD model considered here
~\cite{poon_law_stochastic_pmd_2008,
riccardi_polarization_entanglement_distribution_2021}.

It is convenient to introduce the dimensionless normalized PMD strength

\begin{equation}
\boxed{
\chi_j
=
\sigma_\omega |\Delta\tau_j|.
}
\label{eq:normalized_pmd_parameter}
\end{equation}

The coherence factor can then be written simply as

\begin{equation}
\boxed{
|\kappa_j|
=
e^{-\chi_j^2/2}.
}
\label{eq:gaussian_pmd_coherence_normalized}
\end{equation}

The parameter \(\chi_j\) compares the PMD-induced differential delay with the
spectral scale of the photon. The same DGD therefore produces negligible
polarization decoherence for a sufficiently narrow spectrum and increasingly
strong decoherence as the occupied bandwidth increases.

For the single-arm configuration considered later in the numerical analysis,
with the reference Bell state \( |\Psi^+\rangle \) and one photon undergoing the
fixed-axis PMD transformation while the other arm remains ideal, the spectral
trace produces the same dephased Bell-state family introduced in
Section~\ref{sec:statistical_phase_fluctuations}, with

\begin{equation}
|\mu|
=
|\kappa|
=
e^{-\chi^2/2}.
\label{eq:single_arm_pmd_effective_coherence}
\end{equation}

Using Eq.~\eqref{eq:random_phase_concurrence_tangle}, the concurrence is therefore

\begin{equation}
\boxed{
C(\chi)
=
e^{-\chi^2/2}.
}
\label{eq:single_arm_pmd_concurrence_analytical}
\end{equation}

This result provides the analytical reference for the single-arm PMD numerical
experiment of Chapter~\ref{ch:03_numerical_simulations_results}.

%--------------------------------------------------------
% Subsection 2.7.3 — Fixed-Frequency Limit
%--------------------------------------------------------

\subsection{Fixed-Frequency Limit and Connection to the Residual Phase Model}
\label{subsec:fixed_frequency_birefringence_residual_phase}

At a fixed frequency, the birefringent transformation remains fully unitary. In
the local eigenpolarization basis it can be written, apart from a global phase,
as

\begin{equation}
U_j^{\mathrm{eig}}(\omega_0)
\sim
\begin{pmatrix}
e^{-i\Delta\phi_j/2} & 0\\
0 & e^{+i\Delta\phi_j/2}
\end{pmatrix}.
\label{eq:fixed_frequency_eigenbasis_operator}
\end{equation}

When the relevant basis is aligned with the computational polarization basis,
this reduces to

\begin{equation}
U_j^{\mathrm{eig}}(\omega_0)
\sim
\begin{pmatrix}
1 & 0\\
0 & e^{i\Delta\phi_j}
\end{pmatrix}.
\label{eq:fixed_frequency_reduced_phase_form}
\end{equation}

This is the same operator form introduced in the residual phase model of
Section~\ref{sec:reduced_phase_model}. More generally, polarization compensation
acts on the complete fixed-frequency fiber transformation and leaves, within the
reduced description adopted in this work,

\begin{equation}
U_{\mathrm{fib}}^{\mathrm{comp}}(\omega_0)
\sim
\begin{pmatrix}
1 & 0\\
0 & e^{i\theta}
\end{pmatrix}.
\label{eq:compensated_complete_fiber_reduced_operator}
\end{equation}

The residual phase model can therefore be interpreted as an effective
fixed-frequency limit of the compensated fiber transformation. The additional
frequency dependence introduced by PMD is what allows polarization coherence to
decrease after spectral averaging.


%========================================================
% SECTION 2.8 — Concatenated Birefringent Fiber
%========================================================

\section{Concatenated Birefringent Fiber}
\label{sec:concatenated_birefringent_fiber}

The birefringence of a real fiber generally varies along the propagation
direction. A nonuniform fiber can therefore be represented as a sequence of
locally uniform birefringent segments, each characterized by its own
birefringence axis and differential phase
~\cite{wai_menyuk_random_birefringence_1996,
gordon_kogelnik_pmd_2000}.

Figure~\ref{fig:segmented_fiber_model} illustrates this representation.

\begin{figure}[htbp]
    \centering

    \begin{tikzpicture}[
        >=Latex,
        every node/.style={font=\small}
    ]

        % Input arrow
        \draw[->, thick]
        (-1.0,0)
        --
        (0,0);

        \node[left] at (-1.0,0)
        {\( |\psi_{\mathrm{in}}\rangle \)};

        % Segment rectangles
        \draw (0,-0.65) rectangle (1.8,0.65);
        \draw (1.8,-0.65) rectangle (3.6,0.65);
        \draw (3.6,-0.65) rectangle (5.4,0.65);
        \draw (5.4,-0.65) rectangle (7.2,0.65);
        \draw (7.2,-0.65) rectangle (9.0,0.65);

        % Propagation line
        \draw[->, thick]
        (0,0)
        --
        (9.9,0);

        % Local birefringence axes
        \draw[->, thick]
        (0.90,-0.35)
        --
        (0.90,0.42);

        \draw[->, thick]
        (3.10,0.15)
        --
        (2.10,-0.15);

        \draw[->, thick]
        (4.20,0.32)
        --
        (4.85,-0.42);

        \draw[->, thick]
        (6.45,-0.35)
        --
        (6.05,0.38);

        \draw[->, thick]
        (8.45,0.32)
        --
        (7.70,-0.42);

        % Segment labels
        \node at (0.90,0.95)
        {\(U_{K,1}(\omega)\)};

        \node at (2.70,0.95)
        {\(U_{K,2}(\omega)\)};

        \node at (4.50,0.95)
        {\(\cdots\)};

        \node at (6.30,0.95)
        {\(U_{K,N_K-1}(\omega)\)};

        \node at (8.10,0.95)
        {\(U_{K,N_K}(\omega)\)};

        % Axis labels
        \node at (0.90,-1.05)
        {\(\hat{\mathbf n}_{K,1}\)};

        \node at (2.70,-1.05)
        {\(\hat{\mathbf n}_{K,2}\)};

        \node at (4.50,-1.05)
        {\(\cdots\)};

        \node at (6.30,-1.05)
        {\(\hat{\mathbf n}_{K,N_K-1}\)};

        \node at (8.10,-1.05)
        {\(\hat{\mathbf n}_{K,N_K}\)};

        % Output
        \node[right] at (9.9,0)
        {\( |\psi_{\mathrm{out}}\rangle \)};

    \end{tikzpicture}

    \caption{
        Segmented representation of a nonuniform birefringent fiber. Each
        section is characterized by its own birefringence axis and
        frequency-dependent unitary transformation. The complete fiber
        transformation is obtained from the ordered product of the segment
        operators.
    }
    \label{fig:segmented_fiber_model}
\end{figure}

For arm \(K\), let \(N_K\) denote the number of segments. The complete
frequency-dependent transformation is

\begin{equation}
\boxed{
U_K(\omega_K)
=
U_{K,N_K}(\omega_K)
\cdots
U_{K,2}(\omega_K)
U_{K,1}(\omega_K).
}
\label{eq:concatenated_fiber_operator}
\end{equation}

The ordering is essential because segments with different birefringence axes do
not generally commute,

\begin{equation}
U_{K,j}(\omega)
U_{K,k}(\omega)
\neq
U_{K,k}(\omega)
U_{K,j}(\omega).
\label{eq:noncommuting_birefringent_segments}
\end{equation}

A particular fiber realization is therefore specified by the set of local
birefringence axes, differential phases, and differential delays of its segments.
For a fixed realization and resolved frequency, the complete propagation remains
unitary.

Statistical sampling of the segment parameters can instead be used to represent
different possible spatial birefringence profiles. This statistical description
of the fiber must be distinguished from the reduced polarization mixedness
produced by tracing over unresolved frequency.

%========================================================
% SECTION 2.9 — PMD Acting on Entangled Photon Pairs
%========================================================

\section{PMD Acting on Polarization-Entangled Photon Pairs}
\label{sec:pmd_entangled_photon_pairs}

For the two-arm system, each resolved frequency pair undergoes the local
transformation

\begin{equation}
U_{AB}(\omega_A,\omega_B)
=
U_A(\omega_A)
\otimes
U_B(\omega_B).
\label{eq:pmd_two_arm_operator}
\end{equation}

After propagation and spectral tracing, the observable polarization state is

\begin{equation}
\boxed{
\rho_{AB}^{\mathrm{pol,out}}
=
\iint
d\omega_A\,d\omega_B\,
J(\omega_A,\omega_B)
\left[
U_A(\omega_A)
\otimes
U_B(\omega_B)
\right]
\rho_{AB}^{\mathrm{pol,in}}
\left[
U_A(\omega_A)
\otimes
U_B(\omega_B)
\right]^\dagger.
}
\label{eq:pmd_two_arm_reduced_polarization_state}
\end{equation}

For every individual frequency pair, the transformation remains a product of
local unitary operators. The reduced polarization state can nevertheless become
mixed because the unresolved spectral components undergo different local
transformations before being combined by the spectral trace
~\cite{poon_law_stochastic_pmd_2008,
antonelli_shtaif_brodsky_pmd_2011,
riccardi_polarization_entanglement_distribution_2021}.

The joint spectral distribution also determines how the two local
frequency-dependent transformations are sampled. If

\begin{equation}
J(\omega_A,\omega_B)
=
J_A(\omega_A)
J_B(\omega_B),
\label{eq:factorized_joint_spectrum}
\end{equation}

the effective polarization channel factorizes into independently averaged local
channels,

\begin{equation}
\mathcal{E}_{AB}
=
\mathcal{E}_A
\otimes
\mathcal{E}_B.
\label{eq:factorized_effective_pmd_channel}
\end{equation}

For spectrally correlated photon pairs, this factorization does not generally
hold. Although the physical transformations remain local, the local fiber
responses are averaged according to the common joint spectral distribution of the
photon pair.

The resulting polarization degradation therefore depends jointly on the source
spectrum and on the frequency-dependent transformations of the two fiber arms.

%========================================================
% SECTION 2.10 — Relation Between the Modeling Levels
%========================================================

\section{Relation Between the Modeling Levels}
\label{sec:relation_between_modeling_levels}

The models developed in this chapter retain progressively different amounts of
physical information. The residual phase model isolates coherent polarization
evolution, while statistical phase averaging introduces effective dephasing
through the coherence parameter \(\mu\). The depolarizing models provide
analytical references for isotropic polarization degradation, and the composite
model combines these effects within a polarization-only description.

The frequency-dependent model instead retains the spectral degree of freedom.
Fiber birefringence produces frequency-dependent polarization transformations,
while PMD characterizes their first-order variation through the PSPs and the
differential group delay. For a fixed lossless fiber realization, the complete
polarization--frequency evolution remains unitary, and reduced polarization
decoherence emerges after the unresolved spectral degree of freedom is traced
out.

The resulting hierarchy can be summarized as

\begin{figure}[H]
    \centering

    \begin{tikzpicture}[
        >=Latex,
        every node/.style={font=\small},
        modelbox/.style={
            draw,
            rounded corners,
            align=center,
            text width=0.55\textwidth,
            minimum height=0.72cm,
            inner sep=4pt
        }
    ]

        \node[modelbox] (coh) at (0,0)
        {coherent polarization evolution};

        \node[modelbox] (phase) at (0,-1.3)
        {statistical phase averaging};

        \node[modelbox] (depol) at (0,-2.6)
        {effective depolarization};

        \node[modelbox] (composite) at (0,-3.9)
        {composite polarization channel};

        \node[modelbox] (pmd) at (0,-5.2)
        {frequency-dependent birefringent propagation and PMD};

        \draw[->] (coh) -- (phase);
        \draw[->] (phase) -- (depol);
        \draw[->] (depol) -- (composite);
        \draw[->] (composite) -- (pmd);

    \end{tikzpicture}

    \caption{
        Summary of the modeling hierarchy developed in this chapter, from
        polarization-only effective descriptions to frequency-dependent
        birefringent propagation and PMD.
    }
    \label{fig:complete_modeling_hierarchy}
\end{figure}

The PMD model finally provides the connection between the physical properties of
the source and fiber and the reduced polarization state. The relevant physical
parameters can therefore be collected in the parameter vector
\(\boldsymbol{\xi}\), introduced in
Section~\ref{sec:bbm92_operational_requirements}, while the resulting state and
operational quantities are represented by
\(\boldsymbol{\lambda}(\boldsymbol{\xi})\).

This establishes the physical-to-operational mapping

\begin{equation}
\boxed{
\boldsymbol{\xi}
\longrightarrow
\rho_{AB}^{\mathrm{pol,out}}(\boldsymbol{\xi})
\longrightarrow
\boldsymbol{\lambda}(\boldsymbol{\xi})
}
\label{eq:pmd_physical_operational_mapping}
\end{equation}

that will be investigated numerically in
Chapter~\ref{ch:03_numerical_simulations_results}. In particular, the numerical
analysis will determine how variations of the source and fiber parameters modify
the distributed polarization state and the corresponding region of parameter
space compatible with the adopted BBM92 operational criterion.

%========================================================
% SECTION 2.11 — BBM92 Operating Regions of the Fiber Models
%========================================================

\section{BBM92 Operating Regions of the Fiber Models}
\label{sec:bbm92_operating_regions_fiber_models}

The framework introduced in
Section~\ref{sec:bbm92_operational_requirements} can now be applied to the
specific fiber-channel models developed in this chapter. Each model defines a
different parameter space and therefore a corresponding BBM92 operating region.

For the reference state \( |\Psi^+\rangle \), measurements in the computational
basis are ideally anticorrelated,

\begin{equation}
|\Psi^+\rangle
=
\frac{
|H\rangle_A|V\rangle_B
+
|V\rangle_A|H\rangle_B
}{\sqrt{2}},
\nonumber
\end{equation}

whereas measurements in the diagonal basis are ideally correlated. Indeed,

\begin{equation}
|\Psi^+\rangle
=
\frac{
|D\rangle_A|D\rangle_B
-
|A\rangle_A|A\rangle_B
}{\sqrt{2}}.
\end{equation}

The definition of an erroneous event must therefore follow the expected
correlation pattern in each basis.

In the computational basis, an error corresponds to equal outcomes,
\(HH\) or \(VV\). The corresponding QBER is therefore

\begin{equation}
Q_Z
=
P(HH)+P(VV).
\label{eq:qber_z_probabilities}
\end{equation}

The \(zz\) correlation is

\begin{equation}
T_{zz}
=
P(HH)+P(VV)-P(HV)-P(VH).
\label{eq:tzz_joint_probabilities}
\end{equation}

Since the four joint probabilities sum to unity,

\begin{equation}
P(HV)+P(VH)
=
1-
\left[
P(HH)+P(VV)
\right],
\end{equation}

and consequently

\begin{equation}
T_{zz}
=
2Q_Z-1.
\end{equation}

It follows that

\begin{equation}
\boxed{
Q_Z
=
\frac{1+T_{zz}}{2}.
}
\label{eq:bbm92_qber_z_from_correlation}
\end{equation}

In the diagonal basis, the ideal outcomes are instead equal, \(DD\) or \(AA\).
An error therefore corresponds to the opposite outcomes \(DA\) or \(AD\),

\begin{equation}
Q_X
=
P(DA)+P(AD).
\label{eq:qber_x_probabilities}
\end{equation}

The corresponding \(xx\) correlation is

\begin{equation}
T_{xx}
=
P(DD)+P(AA)-P(DA)-P(AD).
\label{eq:txx_joint_probabilities}
\end{equation}

Using normalization gives

\begin{equation}
T_{xx}
=
1-2Q_X,
\end{equation}

and therefore

\begin{equation}
\boxed{
Q_X
=
\frac{1-T_{xx}}{2}.
}
\label{eq:bbm92_qber_x_from_correlation}
\end{equation}

Thus, for the \( |\Psi^+\rangle \) convention adopted throughout this work, the
basis-dependent QBERs can be obtained directly from the corresponding diagonal
components of the correlation tensor. These relations provide the connection
between the polarization-correlation models developed above and the BBM92
operating regions considered below.

%--------------------------------------------------------
% Subsection 2.11.1 — Residual Phase Model
%--------------------------------------------------------

\subsection{Residual Phase Model}
\label{subsec:bbm92_region_residual_phase}

For the fixed residual-phase model,

\begin{equation}
T_{zz}
=
-1,
\qquad
T_{xx}
=
\cos\theta.
\end{equation}

The corresponding QBERs are therefore

\begin{equation}
Q_Z(\theta)
=
0,
\qquad
Q_X(\theta)
=
\frac{1-\cos\theta}{2}.
\label{eq:bbm92_qber_residual_phase}
\end{equation}

The residual phase does not affect the computational-basis outcomes, but it can
produce errors in the diagonal basis despite the state remaining maximally
entangled.

The corresponding operating region is

\begin{equation}
\boxed{
\mathcal{R}_{\mathrm{BBM92}}^{(\theta)}
=
\left\{
\theta:
\frac{1-\cos\theta}{2}
<
Q_{\mathrm{th}}
\right\}.
}
\label{eq:bbm92_region_residual_phase}
\end{equation}

%--------------------------------------------------------
% Subsection 2.11.2 — Statistical Phase Model
%--------------------------------------------------------

\subsection{Statistical Phase Model}
\label{subsec:bbm92_region_statistical_phase}

For the statistical phase model,

\begin{equation}
T_{zz}
=
-1,
\qquad
T_{xx}
=
\operatorname{Re}(\mu),
\end{equation}

and consequently

\begin{equation}
Q_Z(\mu)
=
0,
\qquad
Q_X(\mu)
=
\frac{
1-\operatorname{Re}(\mu)
}{2}.
\label{eq:bbm92_qber_statistical_phase}
\end{equation}

The BBM92 operating region in the complex coherence-parameter space is therefore

\begin{equation}
\boxed{
\mathcal{R}_{\mathrm{BBM92}}^{(\mu)}
=
\left\{
\mu:
\frac{
1-\operatorname{Re}(\mu)
}{2}
<
Q_{\mathrm{th}}
\right\}.
}
\label{eq:bbm92_region_statistical_phase}
\end{equation}

This result also shows that the operational condition depends on the residual
coherent phase contained in \(\mu\), and not only on its magnitude
\( |\mu|=C \).

%--------------------------------------------------------
% Subsection 2.11.3 — Local Depolarizing Model
%--------------------------------------------------------

\subsection{Local Depolarizing Model}
\label{subsec:bbm92_region_local_depolarization}

For the local two-arm depolarizing model applied to the reference Bell state,

\begin{equation}
T_{zz}
=
-\eta,
\qquad
T_{xx}
=
\eta,
\end{equation}

where

\begin{equation}
\eta
=
(1-p_A)(1-p_B).
\end{equation}

The two measurement bases therefore give the same error rate,

\begin{equation}
Q_Z(\eta)
=
Q_X(\eta)
=
\frac{1-\eta}{2}.
\label{eq:bbm92_qber_local_depolarization}
\end{equation}

The operating region is consequently

\begin{equation}
\boxed{
\mathcal{R}_{\mathrm{BBM92}}^{(\eta)}
=
\left\{
\eta:
\frac{1-\eta}{2}
<
Q_{\mathrm{th}}
\right\}.
}
\label{eq:bbm92_region_local_depolarization}
\end{equation}

Equivalently, in terms of the individual arm parameters,

\begin{equation}
\mathcal{R}_{\mathrm{BBM92}}^{(p_A,p_B)}
=
\left\{
(p_A,p_B):
\frac{
1-(1-p_A)(1-p_B)
}{2}
<
Q_{\mathrm{th}}
\right\}.
\label{eq:bbm92_region_local_depolarization_arm_parameters}
\end{equation}

%--------------------------------------------------------
% Subsection 2.11.4 — Composite Effective Model
%--------------------------------------------------------

\subsection{Composite Effective Model}
\label{subsec:bbm92_region_composite}

For the statistical composite model developed in
Section~\ref{subsec:statistical_composite_model}, the relevant tensor components
are

\begin{equation}
T_{zz}
=
-\eta,
\qquad
T_{xx}
=
\eta\operatorname{Re}(\mu).
\end{equation}

The corresponding error rates are

\begin{equation}
Q_Z(\eta,\mu)
=
\frac{1-\eta}{2},
\qquad
Q_X(\eta,\mu)
=
\frac{
1-\eta\operatorname{Re}(\mu)
}{2}.
\label{eq:bbm92_qber_composite_model}
\end{equation}

The operating region can therefore be expressed as

\begin{equation}
\boxed{
\mathcal{R}_{\mathrm{BBM92}}^{(\eta,\mu)}
=
\left\{
(\eta,\mu):
Q_Z(\eta,\mu)<Q_{\mathrm{th}},
\;
Q_X(\eta,\mu)<Q_{\mathrm{th}}
\right\}.
}
\label{eq:bbm92_region_composite_model}
\end{equation}

This region provides a simple analytical example of an operational boundary in a
multidimensional channel-parameter space.

%--------------------------------------------------------
% Subsection 2.11.5 — Frequency-Dependent PMD Model
%--------------------------------------------------------

\subsection{Frequency-Dependent PMD Model}
\label{subsec:bbm92_region_pmd}

The PMD model replaces the phenomenological parameters of the previous
descriptions with quantities associated with the source spectrum and the physical
fiber propagation.

Let

\begin{equation}
\boldsymbol{\xi}_{\mathrm{PMD}}
=
\{
\sigma_\omega,\,
r,\,
\tau_{\mathrm{DGD},A},\,
\tau_{\mathrm{DGD},B},\,
\hat{\mathbf n}_A,\,
\hat{\mathbf n}_B,\,
\ldots
\}
\label{eq:pmd_physical_parameter_vector}
\end{equation}

denote the relevant set of physical parameters for a particular PMD model. The
reduced polarization state is obtained from
Eq.~\eqref{eq:pmd_two_arm_reduced_polarization_state},

\begin{equation}
\boldsymbol{\xi}_{\mathrm{PMD}}
\longrightarrow
\rho_{AB}^{\mathrm{pol,out}}
\left(
\boldsymbol{\xi}_{\mathrm{PMD}}
\right),
\label{eq:pmd_parameter_to_state_mapping}
\end{equation}

from which \(Q_Z\) and \(Q_X\) are evaluated.

Unlike the effective models above, a universal closed-form operating boundary is
not expected for a general frequency-dependent fiber. The corresponding region is
therefore defined as

\begin{equation}
\boxed{
\mathcal{R}_{\mathrm{BBM92}}^{(\mathrm{PMD})}
=
\left\{
\boldsymbol{\xi}_{\mathrm{PMD}}:
Q_Z(\boldsymbol{\xi}_{\mathrm{PMD}})<Q_{\mathrm{th}},
\;
Q_X(\boldsymbol{\xi}_{\mathrm{PMD}})<Q_{\mathrm{th}}
\right\}.
}
\label{eq:bbm92_region_pmd}
\end{equation}

This formulation converts the PMD model into an operational problem in physical
parameter space. The numerical simulations of
Chapter~\ref{ch:03_numerical_simulations_results} are used to evaluate these
regions and determine how their boundaries depend on the source, fiber, and
compensation parameters.

%========================================================
% END OF CHAPTER 2
%========================================================